\documentclass[11pt,a4paper]{article}

\usepackage[T1]{fontenc}
\usepackage[utf8]{inputenc}
\usepackage[french]{babel}
\usepackage[margin=2.4cm]{geometry}
\usepackage{booktabs}
\usepackage{longtable}
\usepackage{array}
\usepackage{xcolor}
\usepackage{enumitem}
\usepackage{fancyhdr}
\usepackage[hidelinks]{hyperref}

\definecolor{crit}{HTML}{A13D3D}
\definecolor{warn}{HTML}{B45309}
\definecolor{ok}{HTML}{2F6F4F}
\definecolor{grey}{HTML}{5A6472}

\newcommand{\ko}{\textcolor{crit}{\textbf{$\times$}}}
\newcommand{\att}{\textcolor{warn}{\textbf{!}}}
\newcommand{\oui}{\textcolor{ok}{\textbf{$\checkmark$}}}
\newcommand{\cle}[1]{\texttt{\small #1}}

\pagestyle{fancy}
\fancyhf{}
\lhead{\small\textsc{Audit des références} — \textit{Language Doesn't Matter}}
\rhead{\small\thepage}
\renewcommand{\headrulewidth}{0.4pt}

\setlist{itemsep=2pt,topsep=4pt}

\title{\vspace{-1.5cm}\textbf{Audit des références}\\[2mm]
\large \textit{Language Doesn't Matter} --- commit \texttt{361b2c4}}
\author{}
\date{10 septembre 2026}

\begin{document}
\maketitle
\thispagestyle{fancy}
\vspace{-1cm}

%======================================================================
\section*{Ce qu'il faut savoir en une page}
\addcontentsline{toc}{section}{Résumé}

Les 32 références citées par le manuscrit ont été vérifiées auprès des
registraires (Crossref, DataCite, PubMed, ACL Anthology). \textbf{Aucune n'est
inventée} : toutes existent et les métadonnées du fichier \cle{references.bib}
sont exactes.

Le problème est ailleurs. \textbf{Trois références réelles sont citées pour des
affirmations qu'elles ne soutiennent pas}, et une quatrième porte une
affirmation factuellement inexacte. Une cinquième erreur a été introduite le
9~septembre lors du dégraissage du manuscrit et a déjà été corrigée.

\begin{center}
\begin{tabular}{@{}p{3.6cm}p{11.5cm}@{}}
\toprule
\textbf{Référence} & \textbf{Problème} \\
\midrule
\cle{Jiang2020How} &
Article LPAQA sur le sondage de connaissances factuelles \emph{en anglais}
(« Obama is a \_\_ by profession »). Cité \textbf{trois fois} pour fonder
l'avantage du prompt anglais en contexte translinguistique. \\[2mm]
\cle{Plagwitz2024Zero} &
Reconnaissance d'entités nommées sur des textes cliniques \emph{allemands}.
Cité pour étayer l'analyse de sentiment en italien et en russe. \\[2mm]
\cle{Vileikyte2024Sentiment} &
Fine-tuning de BERT et T5 sur des avis lituaniens. Cité pour un « avantage de
la langue locale dans le prompt », question que l'article ne pose pas. \\[2mm]
\cle{Venerito2024Large} &
Les entretiens italiens sont \emph{traduits en anglais} avant analyse. Cité
comme démonstration de performance « en italien ». \\
\bottomrule
\end{tabular}
\end{center}

\vspace{3mm}
Par ailleurs, l'affirmation chiffrée la plus vérifiable du manuscrit --- «~often
above 85-95\% accuracy across Indo-European, Semitic, and Asian languages~» ---
\textbf{n'a pas pu être confirmée} : les éditeurs bloquent le téléchargement
automatisé et les résumés ne contiennent pas ces chiffres. La taxonomie qui
l'accompagne est en outre incohérente (voir §\ref{sec:taxo}).

%======================================================================
\section{Méthode et niveaux de preuve}

Trois niveaux de preuve ont été employés. \textbf{Ils ne sont pas
équivalents}, et la distinction gouverne tout ce document.

\begin{center}
\begin{tabular}{@{}llp{9.4cm}@{}}
\toprule
\textbf{Niveau} & \textbf{Couverture} & \textbf{Ce qu'il permet d'affirmer} \\
\midrule
Métadonnées & 32 / 32 &
L'article existe ; auteurs, titre, année et DOI sont exacts. Source :
API du registraire. \\[1.5mm]
Résumé & 21 / 32 &
Le sujet de l'article, tel que ses auteurs le décrivent. Suffit à détecter
un contresens sur l'objet, pas sur un résultat chiffré. \\[1.5mm]
Texte intégral & 0 / 32 &
Rien n'a été lu en entier. Aucune affirmation numérique n'est vérifiée. \\
\bottomrule
\end{tabular}
\end{center}

\vspace{2mm}
\noindent\textbf{Limite à retenir.} Un audit fondé sur les résumés attrape les
erreurs d'objet --- citer un article de reconnaissance d'entités pour de
l'analyse de sentiment --- mais pas les erreurs de degré. Si un article rapporte
78\,\% d'exactitude et que le manuscrit écrit 85\,\%, cet audit ne le voit pas.
C'est pourquoi la liste~A existe.

\medskip
\noindent\textbf{Un indice de provenance.} Sur les 101 entrées du fichier
\cle{references.bib}, 62 portent un champ \cle{file} pointant vers une
bibliothèque Zotero (\cle{/home/ral/Zotero/storage/...}) absente de ce dépôt.
Parmi les 32~références citées, \textbf{23 ont un PDF attaché et 9 n'en ont
pas}. Ces neuf ont vraisemblablement été ajoutées par import de métadonnées,
sans que le texte intégral soit consulté. Elles figurent toutes en liste~A ou~B.

%======================================================================
\section{Liste A --- à lire intégralement}
\label{sec:A}

Dix références. Soit un problème y est avéré, soit elles portent une affirmation
chiffrée que seul le texte intégral permet de vérifier. \textbf{Aucune ne devrait
rester dans le manuscrit sans lecture.}

\subsection*{A.1 --- Erreurs avérées (à corriger, pas seulement à lire)}

\begin{description}[leftmargin=0pt,style=unboxed]

\item[\ko~\cle{Jiang2020How}] --- Jiang, Xu, Araki \& Neubig (2020),
\emph{How Can We Know What Language Models Know?}, TACL.\\
\textit{Le manuscrit affirme} : «~a growing body of research reports that
prompting in English improves performance on non-English tasks~».\\
\textit{L'article traite de} : la génération automatique de prompts pour
extraire des connaissances factuelles d'un modèle, en anglais. Le résumé ne
mentionne ni multilinguisme ni langue de prompt.\\
\textit{Action} : retirer des trois emplacements (Introduction, Revue,
Discussion). \textbf{Hypothèse à vérifier} : les mêmes auteurs ont publié en
2020 \emph{X-FACTR: Multilingual Factual Knowledge Retrieval from Pretrained
Language Models}, qui est multilingue. Une confusion d'entrée Zotero est
plausible ; si X-FACTR soutient l'affirmation, la correction se réduit à un
changement de clé.

\item[\ko~\cle{Plagwitz2024Zero}] --- Plagwitz et al. (2024), \emph{Zero-Shot
LLMs for Named Entity Recognition: Targeting Cardiac Function Indicators}.\\
\textit{Le manuscrit affirme} : «~Preliminary work with Mistral-7B, LLaMA2, and
instruction-tuned models suggests promise in Italian, Russian, and cross-lingual
settings~».\\
\textit{L'article traite de} : extraction de paramètres médicaux depuis des
textes cliniques \emph{allemands}, par reconnaissance d'entités nommées. Ni
analyse de sentiment, ni italien, ni russe.\\
\textit{Action} : retirer de cette grappe. C'est aussi la seule des quatre
références à mentionner Mistral --- retirer la référence oblige donc à retirer
« Mistral-7B » de l'affirmation, ou à trouver une autre source.

\item[\ko~\cle{Vileikyte2024Sentiment}] --- Vileikytė et al. (2024),
\emph{Sentiment Analysis of Lithuanian Online Reviews Using Large Language
Models}.\\
\textit{Le manuscrit affirme} : «~some studies find language-specific prompting
advantages~».\\
\textit{L'article traite de} : le fine-tuning de BERT et T5 sur des avis
lituaniens, avec la conclusion que les modèles affinés battent GPT-4 (80,7\,\%
et 89,6\,\% d'exactitude). C'est une comparaison \emph{fine-tuning contre
généraliste}, pas une comparaison de langues de prompt.\\
\textit{Action} : retirer de cette grappe. La phrase « evidence is mixed »
repose alors sur \cle{Hamalainen2022Video} seul --- à vérifier d'urgence
(liste~B), faute de quoi l'affirmation d'un débat n'a plus de support.

\item[\att~\cle{Venerito2024Large}] --- Venerito et al. (2024), \emph{Large
Language Model-Driven Sentiment Analysis for Facilitating Fibromyalgia
Diagnosis}, RMD Open.\\
\textit{Le manuscrit affirme} : « promise in Italian ».\\
\textit{L'article précise} : «~Transcribed responses to questions on pain and
sleep were \emph{machine translated to English} and analysed by the LLM
Mistral-7B-Instruct-v0.2~». L'analyse ne porte donc pas sur de l'italien.\\
\textit{Action} : corriger l'affirmation. \textbf{À noter} : ce cas est en
réalité \emph{intéressant} pour votre article --- c'est une condition
EN$\rightarrow$EN appliquée en clinique. Bien reformulé, il devient un appui
plutôt qu'une erreur.

\end{description}

\subsection*{A.2 --- L'affirmation chiffrée non vérifiée}

Cinq références étayent «~often above 85-95\% accuracy across Indo-European,
Semitic, and Asian languages~». Aucun chiffre n'a pu être confirmé.

\begin{center}
\begin{tabular}{@{}llll@{}}
\toprule
\textbf{Référence} & \textbf{Langue} & \textbf{Famille} & \textbf{PDF} \\
\midrule
\cle{Anirban2021Sentiment}      & bengali    & indo-européenne & \textbf{sans DOI} \\
\cle{Michailidis2024Comparative} & grec       & indo-européenne & MDPI, 403 \\
\cle{Rafeef2023Evaluation}       & arabe      & sémitique       & accessible \\
\cle{ElJundi2019hULMonA}         & arabe      & sémitique       & ACL, lien mort \\
\cle{Ahmadian2024Hybrid}         & indonésien & austronésienne  & Wiley, 403 \\
\bottomrule
\end{tabular}
\end{center}

\medskip
\noindent\textbf{Deux problèmes distincts.}

\begin{enumerate}
\item \textbf{Le chiffre.} Il faut ouvrir les cinq articles et relever
l'exactitude rapportée. Si l'un d'eux se situe à 70\,\%, l'intervalle « 85-95 »
est faux. C'est la vérification la plus facile à faire pour un réviseur, donc
la plus risquée à négliger.

\item \textbf{Une référence détonne.} \cle{Rafeef2023Evaluation} porte sur
l'évaluation d'hôtels à partir d'avis en arabe, publiée dans une revue de
technologies mobiles interactives. C'est une application sectorielle, appui
faible pour une affirmation générale sur la performance translinguistique.
Envisager de la retirer.
\end{enumerate}

\subsection*{A.3 --- Un soutien réel, mais sur un autre mécanisme}

\begin{description}[leftmargin=0pt,style=unboxed]
\item[\att~\cle{Muennighoff2023Crosslingual}] --- Muennighoff et al. (2023),
\emph{Crosslingual Generalization through Multitask Finetuning}, ACL.\\
Le résumé dit : «~finetuning large multilingual language models on English tasks
with English prompts allows for task generalization to non-English languages~»
et «~Finetuning on multilingual tasks with English prompts further improves
performance~».\\
\textit{Nuance} : il s'agit d'un effet du \textbf{fine-tuning} avec prompts
anglais, non du prompting à l'inférence --- ce que votre étude mesure. Les deux
mécanismes sont distincts.\\
\textit{Action} : conserver, mais préciser la distinction dans le texte. C'est à
votre avantage : cela montre que la question posée à l'inférence reste ouverte,
ce qui renforce la justification de l'étude.
\end{description}

%======================================================================
\section{Liste B --- à vérifier sur un point précis}

Quatorze références. Le résumé est cohérent avec l'usage qu'en fait le
manuscrit, ou n'a pas pu être récupéré. Une lecture ciblée suffit : résumé,
tableau de résultats, conclusion.

\begin{center}
\small
\begin{longtable}{@{}p{3.4cm}p{7.6cm}p{3.4cm}@{}}
\toprule
\textbf{Référence} & \textbf{Point à vérifier} & \textbf{État} \\
\midrule
\endhead
\cle{Hamalainen2022Video} &
Trouve-t-il vraiment un avantage de la langue locale ? Devient le \emph{seul}
support de « evidence is mixed » si \cle{Vileikyte} est retirée. &
\textbf{prioritaire} \\
\cle{Lin2022Few} &
Contient-il des expériences sur la \emph{langue du prompt} ? Le résumé porte sur
la capacité multilingue, pas sur le prompting. & résumé lu \\
\cle{Nie2024Decomposed} &
Porte sur l'étiquetage de séquences et le maintien de gabarits de sortie.
Soutient-il l'avantage du prompt anglais ? & résumé lu \\
\cle{Pan2023Evaluation} &
Cité pour le transfert translinguistique ; le résumé porte sur la prédiction
boursière à partir de données financières. & résumé lu \\
\cle{Kumar2021Sentiment} &
Cité pour le transfert translinguistique ; résumé centré sur les médias
sociaux. & résumé lu \\
\cle{Rusnachenko2024Large} &
Russe et LLM affinés par instruction : concordance apparente. Vérification
légère. & résumé lu \\
\cle{Buscemi2024ChatGPT} &
LLaMA2 et 10 langues : concordance apparente. Vérification légère. &
résumé lu \\
\cle{Tela2024Transferring} &
Langues peu dotées : concordance apparente. & résumé lu \\
\cle{Salahudeen2023HausaNLP} &
SemEval-2023 tâche 12, langues africaines : concordance apparente. &
résumé lu \\
\cle{Nazir2025Leveraging} & Aucun résumé récupéré. & \textbf{non vérifié} \\
\cle{Kotelnikova2023Cross} & Aucun résumé récupéré. & \textbf{non vérifié} \\
\cle{Francesco2021XLM} & Aucun DOI, aucun résumé. & \textbf{non vérifié} \\
\cle{Mikhail2024Multilingual} & Aucun DOI, aucun résumé. & \textbf{non vérifié} \\
\cle{Priban2024comparative} & Aucun DOI, aucun résumé. Citée deux fois. &
\textbf{non vérifié} \\
\bottomrule
\end{longtable}
\end{center}

\noindent Les cinq marquées « non vérifié » n'ont \emph{aucune} preuve au-delà
de leur présence dans le fichier \cle{.bib}. Quatre d'entre elles n'ont pas de
DOI, ce qui empêche même la vérification d'existence. \textbf{Ce sont les
premières à traiter dans cette liste.}

%======================================================================
\section{Liste C --- utilisables sans relecture}

Huit références. Le résumé confirme la concordance avec l'usage qu'en fait le
manuscrit, et l'affirmation portée est consensuelle. Le risque est faible et le
coût d'une erreur le serait aussi.

\begin{center}
\begin{tabular}{@{}p{3.6cm}p{11.2cm}@{}}
\toprule
\textbf{Référence} & \textbf{Affirmation portée} \\
\midrule
\cle{touvron\_etal23a} & Les modèles sont entraînés majoritairement sur des
données anglaises. Article Llama, non contesté. \\
\cle{Laurer2023Less} & L'apprentissage supervisé exige des milliers
d'annotations. Résumé : concordance exacte. \\
\cle{Chalkidis2020LEGAL} & L'adaptation de BERT à un domaine exige de
l'expertise. Résumé : concordance. \\
\cle{Strubell2019Energy} & Coût computationnel de l'entraînement. Référence
standard. \\
\cle{Perron2024Demystifying} & Les API rendent les LLM accessibles aux
chercheurs en sciences sociales. Résumé : concordance exacte. \\
\cle{Ghosh2023Navigating} & Même affirmation. Concordance plus lâche
(l'article porte sur GPT-4 en finance et en politique) mais l'enjeu est
faible. \\
\cle{duval\_petry16} & Source du dictionnaire français \cle{frlsd}.
Référence méthodologique standard. \\
\cle{young\_andsoroka12} & Source du \emph{Lexicoder Sentiment Dictionary}.
Référence méthodologique standard. \\
\bottomrule
\end{tabular}
\end{center}

\medskip
\noindent\textbf{Une correction matérielle} : \cle{duval\_petry16} et
\cle{young\_andsoroka12} portent tous deux un champ \cle{date = 2025-04-02},
qui est une date d'import Zotero et non de publication. À corriger dans le
fichier \cle{.bib} : les deux articles datent respectivement de 2016 et 2012.

%======================================================================
\section{Liste D --- les sept références ajoutées le 10 septembre}
\label{sec:D}

Ces sept références ont été ajoutées au fichier \cle{references.bib} pour
combler l'absence de travaux 2025--2026, mais \textbf{ne sont encore citées
nulle part dans le manuscrit}. Elles ne peuvent donc pas encore être mal
citées --- mais elles le seront dès leur intégration si elles ne sont pas lues
d'abord.

\medskip
\noindent\textbf{Leur niveau de preuve est plus faible que celui des 32 autres.}
Pour cinq d'entre elles, la seule information de contenu dont on dispose est un
extrait de moteur de recherche --- une synthèse automatique qui agrège plusieurs
sources sans les distinguer. C'est précisément le mode de défaillance qui
produit une citation mal fondée.

\begin{center}
\small
\begin{tabular}{@{}p{3.3cm}p{5.4cm}p{3.1cm}p{2.6cm}@{}}
\toprule
\textbf{Référence} & \textbf{Ce qu'elle doit étayer} & \textbf{Preuve
disponible} & \textbf{Priorité} \\
\midrule
\cle{richie\_etal22} &
L'accord intercodeurs n'est pas une borne supérieure stricte de la performance
machine. &
Résumé du PDF, lu par un outil automatique &
\textbf{critique} \\[1.5mm]
\cle{nguyen\_etal26} &
La langue du prompt n'est pas neutre : elle modifie la forme de la sortie sans
altérer le sens. &
Résumé de la page arXiv &
\textbf{critique} \\[1.5mm]
\cle{gupta\_etal25} &
Traduire \emph{une partie} du prompt bat le tout-anglais comme le
tout-langue-source. &
Extrait de recherche seulement &
haute \\[1.5mm]
\cle{beauchemin\_etal25} &
Banc d'essai français ; écarts entre modèles à poids ouverts et fermés. &
Extrait de recherche seulement &
haute \\[1.5mm]
\cle{kastrati\_etal26} &
Analyse de sentiment translinguistique par familles de langues. &
Extrait de recherche seulement &
moyenne \\[1.5mm]
\cle{langlais\_etal26} &
Sentiment en français ; modèles spécialisés contre gros LLM. &
Extrait de recherche seulement &
moyenne \\[1.5mm]
\cle{boguslav\_cohen17} &
Antécédent de \cle{richie\_etal22}. &
Métadonnées seulement &
basse \\
\bottomrule
\end{tabular}
\end{center}

\subsection*{Les deux critiques, et pourquoi}

\begin{description}[leftmargin=0pt,style=unboxed]

\item[\att~\cle{richie\_etal22}] --- Richie, Grover \& Tsui (2022), \emph{Inter-annotator
agreement is not the ceiling of machine learning performance}, BioNLP.\\
Cette référence est destinée à étayer une \textbf{reformulation de votre
résultat le plus fort}. Le calcul du plafond humain ($r = 0{,}748$, sept
combinaisons modèle-condition au-dessus) est exact, mais la littérature conteste
que l'accord intercodeurs constitue une borne supérieure. La formulation
« les modèles dépassent la performance humaine » devient donc indéfendable, au
profit de « les modèles atteignent le niveau d'accord que les codeurs humains
obtiennent entre eux ».\\
\textbf{Si cet article ne dit pas ce qui lui est attribué ici, cette
reformulation perd son appui.} C'est la lecture la plus rentable de toute la
liste.

\item[\att~\cle{nguyen\_etal26}] --- Nguyen et al. (2026), \emph{Investigating the
Influence of Prompt and Response Languages on LLM Content Generation}.\\
C'est le travail le plus directement concurrent du vôtre : il conclut que « la
langue du prompt n'est pas neutre », ce qui contredit votre résultat en
apparence. La distinction qui réconcilie les deux --- ils mesurent la
\emph{forme} de la sortie (longueur, réalisation lexicale), vous mesurez une
note numérique, donc le \emph{sens} --- repose sur un résumé, pas sur une
lecture. Si elle ne tient pas, votre positionnement par rapport à ce travail
doit être repensé.

\end{description}

\subsection*{PDF --- six sur sept librement accessibles}

\begin{center}
\begin{tabular}{@{}p{3.6cm}p{11.2cm}@{}}
\toprule
\textbf{Référence} & \textbf{Adresse} \\
\midrule
\cle{richie\_etal22} & \url{https://aclanthology.org/2022.bionlp-1.26.pdf} \\
\cle{nguyen\_etal26} & \url{https://arxiv.org/pdf/2608.26186} \\
\cle{gupta\_etal25} & \url{https://arxiv.org/pdf/2507.22923} \\
\cle{beauchemin\_etal25} & \url{https://arxiv.org/pdf/2510.05046} \\
\cle{langlais\_etal26} & \url{https://arxiv.org/pdf/2604.18226} \\
\cle{kastrati\_etal26} & \url{https://www.frontiersin.org/journals/artificial-intelligence/articles/10.3389/frai.2026.1854873/pdf} \\
\midrule
\cle{boguslav\_cohen17} & IOS Press --- erreur 403. Accès institutionnel requis,
ou à abandonner : \cle{richie\_etal22} couvre le même argument plus récemment. \\
\bottomrule
\end{tabular}
\end{center}

\subsection*{Une remarque sur \cle{beauchemin\_etal25}}

Le banc d'essai COLE est signé par David Beauchemin et Richard Khoury, du
département d'informatique et de génie logiciel de l'\textbf{Université Laval}.
Il porte sur l'évaluation du français et documente les écarts entre modèles à
poids ouverts et fermés --- c'est-à-dire votre question de recherche n\textsuperscript{o}~2,
traitée sur le même campus. Ne pas le citer serait remarqué. Une conversation
avec ces collègues avant soumission pourrait par ailleurs être utile.

%======================================================================
\section{La taxonomie des familles linguistiques}
\label{sec:taxo}

Le manuscrit écrit «~across Indo-European, Semitic, and \emph{Asian}
languages~». L'énumération mêle deux familles linguistiques et une catégorie
géographique, ce qui la rend incohérente :

\begin{itemize}
\item le \textbf{bengali} (\cle{Anirban2021}) est indo-européen \emph{et}
asiatique ;
\item l'\textbf{indonésien} (\cle{Ahmadian2024}) est austronésien, pas
« asiatique » au sens familial ;
\item aucune langue \textbf{sino-tibétaine, japonique ou coréenne} n'est
représentée --- celles qu'un lecteur associe spontanément au mot « asiatiques ».
\end{itemize}

\noindent\textbf{Correction proposée} : «~across Indo-European, Semitic, and
Austronesian languages~», ou plus simplement une énumération par langue
(« en grec, bengali, arabe et indonésien »), qui est à la fois exacte et plus
informative.

%======================================================================
\section{PDF à récupérer}

\subsection*{Librement accessibles --- à télécharger}

\begin{center}
\begin{tabular}{@{}p{4.4cm}p{9.6cm}@{}}
\toprule
\textbf{Référence} & \textbf{Adresse} \\
\midrule
\cle{Vileikyte2024Sentiment} & \url{https://arxiv.org/pdf/2407.19914} \\
\cle{Muennighoff2023Cross...} & \url{https://aclanthology.org/2023.acl-long.891.pdf} \\
\cle{Lin2022Few} & \url{https://aclanthology.org/2022.emnlp-main.616.pdf} \\
\cle{Nie2024Decomposed} & \url{https://arxiv.org/pdf/2402.18397} \\
\cle{Rafeef2023Evaluation} & \url{https://doi.org/10.3991/ijim.v17i09.38755} \\
\bottomrule
\end{tabular}
\end{center}

\subsection*{Bloqués pour un accès automatisé --- accès institutionnel requis}

Ces éditeurs renvoient une erreur 403 aux requêtes automatisées. \textbf{Ils
sont normalement accessibles depuis le réseau de l'Université Laval}, ou via
le proxy de la bibliothèque.

\begin{center}
\begin{tabular}{@{}p{4.4cm}p{4.6cm}p{4.8cm}@{}}
\toprule
\textbf{Référence} & \textbf{Éditeur} & \textbf{Priorité} \\
\midrule
\cle{Jiang2020How} & MIT Press (TACL) & \textbf{critique} \\
\cle{Plagwitz2024Zero} & IOS Press & \textbf{critique} \\
\cle{Venerito2024Large} & BMJ (RMD Open) & \textbf{critique} \\
\cle{Michailidis2024Comparative} & MDPI & haute (chiffre) \\
\cle{Ahmadian2024Hybrid} & Wiley & haute (chiffre) \\
\cle{Hamalainen2022Video} & ACM & haute \\
\cle{Pan2023Evaluation} & PeerJ & moyenne \\
\cle{Kumar2021Sentiment} & ACM & moyenne \\
\bottomrule
\end{tabular}
\end{center}

\subsection*{Problèmes d'accès particuliers}

\begin{itemize}
\item \cle{ElJundi2019hULMonA} --- le DOI ACL (\cle{10.18653/v1/w19-4608})
renvoie une erreur 404 sur l'anthologie. À rechercher manuellement dans les
actes du \emph{Fourth Arabic Natural Language Processing Workshop} (2019).
\item \cle{Anirban2021Sentiment} --- \textbf{aucun DOI dans le fichier
\cle{.bib}}. Impossible de vérifier même son existence. À retrouver et
compléter, ou à retirer.
\item \cle{Francesco2021XLM}, \cle{Mikhail2024Multilingual},
\cle{Priban2024comparative}, \cle{touvron\_etal23a} --- également sans DOI.
Les trois premières sont non vérifiées ; la quatrième est identifiable
(article Llama).
\end{itemize}

\subsection*{La solution durable}

Vingt-trois des trente-deux références citées ont déjà un PDF dans la
bibliothèque Zotero de \cle{ral}. \textbf{Un export Zotero avec pièces jointes,
ou un groupe Zotero partagé}, réglerait la majeure partie du problème d'accès
d'un coup --- et mettrait tous les coauteurs sur la même base pour la suite.
C'est la démarche à privilégier avant d'aller chercher les PDF un par un.

%======================================================================
\section{Les 69 entrées non citées}

Le fichier \cle{references.bib} contient 101 entrées dont 32 seulement sont
citées. Les 69 restantes ne posent aucun problème en l'état --- elles
n'apparaissent pas dans le document rendu.

Elles méritent toutefois un examen avant soumission, pour deux raisons :
certaines couvrent peut-être déjà les lacunes de la revue de littérature
identifiées par ailleurs (absence de travaux 2025--2026), et leur présence
suggère qu'une recherche plus large a été amorcée sans être intégrée au texte.

%======================================================================
\section{Ce que cet audit ne couvre pas}

Par honnêteté sur sa portée :

\begin{itemize}
\item \textbf{Aucun article n'a été lu en entier.} Les verdicts reposent sur les
résumés rédigés par les auteurs, récupérés auprès des registraires.
\item \textbf{Aucune affirmation numérique n'est vérifiée.} L'intervalle
« 85-95\,\% » reste ouvert, ainsi que toute autre valeur chiffrée attribuée à
une source.
\item \textbf{Onze résumés sur trente-deux n'ont pas pu être récupérés.} Pour
ces références, la vérification s'arrête aux métadonnées.
\item \textbf{Les sept références ajoutées le 10 septembre}
(liste~D) reposent sur une preuve encore plus mince : cinq d'entre elles ne sont
étayées que par un extrait de moteur de recherche, et les deux autres par un
résumé produit automatiquement. Leurs métadonnées sont vérifiées ; leur contenu
ne l'est pas.
\item \textbf{L'adéquation entre l'affirmation et la source} a été évaluée sur
l'objet de l'article, non sur la finesse de ses résultats. Un article peut
porter sur le bon sujet et ne pas soutenir la nuance exacte qu'on lui fait
porter.
\end{itemize}

\vspace{4mm}
\noindent\rule{\linewidth}{0.4pt}
\vspace{1mm}

\noindent\small\textcolor{grey}{Audit produit le 10 septembre 2026. Métadonnées
vérifiées auprès de Crossref, DataCite, PubMed et ACL Anthology. Les citations
du manuscrit ont été extraites de \cle{docs/pub/pub.qmd} au commit
\cle{361b2c4}. Une cinquième erreur --- des citations laissées attachées à une
phrase réécrite le 9 septembre --- a été corrigée avant la production de ce
document.}

\end{document}
